\subsection{Applying the convergence estimates to the selected histories}
\label{sec:selected-history-limits}
The preceding convergence arguments concern particular arrays, paths and
stationary branches. We now identify these objects along one construction.
This identification is needed because equality of the actions at one
period would not identify their coefficient arrays, and because a
cutoff-dependent change of the top boundary need not preserve a limit
unless it eventually leaves each fixed scale unchanged.

\subsubsection{One simultaneous stationary construction}
Fix the scale schedule, clipping data and coordinate conventions. Denote
the terminal cutoff by $N$ and the scale by $m\le N$. The finite
stationary construction of Subsection~\ref{sec:path-attachment-producer}
gives paths $R_{N,m}$, scalar corrections $\alpha_{N,m}$ and stationary
maps $\rho_{N,m}^{p,F}$, where $p$ is a period and $F$ an allowed forest.
All these objects are chosen together. Their one-step identities use
$R_{N,m+1}$ as input and $R_{N,m}$ as output. The path bounds hold at
the prescribed scale widths and allowances.

For the convergence assertion we use the strengthened version of this
construction. At scale $m$ it fixes three coefficient Hilbert spaces,
with encoding widths $\delta_{m,2},\delta_{m,3},\delta_{m,4}$ strictly
larger than the paid width $w_m$. The nonlinear path terms are encoded
in their canonical selected subspaces. It also supplies the following
uniformity: every admissible replacement of the input path with the
same attached row and with all path and loop allowances enlarged by
$64$ has an output in these same selected spaces, at the ordinary
output allowances. The output has the same classical and thermal
matching identities on every period and forest. Finally, the underlying
free horizon converges at each scale in all of its retained components,
and the root rows have an eventual bound at some width $v_m>w_m$.
These are the extended-domain, canonical-readout and retained-horizon
hypotheses used in Subsection~\ref{sec:background-path-convergence}.
They belong to the one-step analytic construction; neither finite
backward induction nor the argument below supplies them.

Here is why its coefficient representation is the one needed in that
section. Select the three coefficient arrays of each path once, at
$(N,m)$, and use that same selection in every occurrence of the path.
For $m+2<N$, the input arrays of the step at $m$ are then exactly the
output arrays of the step at $m+1$. The row at either occurrence is
also the same attached root row. Thus the input/output equality holds
in the coefficient spaces, before taking any limit. The one-step
identities and bounds are available for all these chosen arrays.
For a replacement of one array the other arrays remain fixed; their
closed linear constraints and their separate budgets are preserved.
The enlarged input domain therefore supplies exactly the additional
steps used in the contraction argument.

Apply Subsection~\ref{sec:background-path-convergence} to these data.
At every fixed scale there are one limiting path $R_m$ and one
scalar $\alpha_m$ such that every coefficient of $R_{N,m}$ converges
to the corresponding coefficient of $R_m$, and
$\alpha_{N,m}\to\alpha_m$. The limiting path has the original paid
bounds. For every fixed period, forest and derivative order $r$,
\begin{equation}\label{eq:selected-stationary-limit}
 D^r\rho_{N,m}^{p,F}(0)\longrightarrow D^r\rho_m^{p,F}(0).
\end{equation}
The limiting stationary germ is determined by $R_{m+1}$, its retained
quadratic data and its forest. Its Hessian is positive with the same
lower bound. Consequently no choice of a new subsequence for a forest
or derivative order is involved in \eqref{eq:selected-stationary-limit}.
The scalar limit is the one obtained from these same paths and thermal
identities.

\subsubsection{A single width for all four coefficient families}
The common dictionaries are part of the conclusion just obtained.
Write $r_{N,m}(a)$ and $r_m(a)$ for row coefficients on the canonical
row dictionary, and $c_{N,m,q}(a,\nu)$ and $c_{m,q}(a,\nu)$ for the
actual, exponentially encoded coefficients of the atoms of degree
$q=2,3,4$. Their supports, lengths $\ell(a)$ and color multi-indices
$\nu$ are independent of $N$ on a tail. Choose
\[
 w_m<\widetilde w_m<
 \min\{v_m,\delta_{m,2},\delta_{m,3},\delta_{m,4}\}.
\]
Then, for every $\varepsilon>0$, all sufficiently large $N$ satisfy
\begin{align}
 \sum_a |r_{N,m}(a)-r_m(a)|e^{\widetilde w_m\ell(a)}
       &\le\varepsilon,\label{eq:selected-row-width}\\
 \sum_{a,\nu}|c_{N,m,q}(a,\nu)-c_{m,q}(a,\nu)|
                         e^{\widetilde w_m\ell(a)}
       &\le\varepsilon\qquad(q=2,3,4).
 \label{eq:selected-atom-width}
\end{align}
For the row statement one may either sum the finitely many plaquette
directions or impose the estimate separately on each; shrinking the
requested $\varepsilon$ by their number gives the displayed version.

For completeness, the two tail bounds that permit this simultaneous
choice are different. If the row budget at $v_m$ is $B_m$, its
coefficient difference at the smaller width is bounded by
\[2B_m e^{-(v_m-\widetilde w_m)\ell(a)}.\]
For an atom, remove its
encoding factor $e^{-\delta_{m,q}\ell(a)}$. The resulting coefficient
arrays have a common Hilbert bound $C_{m,q}$, so their individual
differences have absolute value at most $2C_{m,q}$. Their weighted
majorant is therefore
$2C_{m,q}e^{-(\delta_{m,q}-\widetilde w_m)\ell(a)}$.
Both majorants are summable on the canonical dictionaries by the
geometric sums in Subsection~\ref{sec:fixed-rung-convergence}. The
coordinate limits and a finite-head, summable-tail argument prove
\eqref{eq:selected-row-width}--\eqref{eq:selected-atom-width}.
There are only four families and finitely many row directions, so
their eventual estimates hold on one tail. This gives one path
limit with one strictly larger width, as required in the local
integral comparison.

\subsubsection{Changing the terminal normalization}
The final Wilson normalization replaces the top path at $m=N$.
Let $\widetilde R_{N,m}$, $\widetilde\alpha_{N,m}$ and
$\widetilde\rho_{N,m}^{p,F}$ denote the resulting data. The change
of coordinates retains the exact identities
\begin{align*}
 \widetilde R_{N,m}&=R_{N,m}&& (m<N),\\
 \widetilde\alpha_{N,m}&=\alpha_{N,m}&&\text{for all }N,m,\\
 \widetilde\rho_{N,m}^{p,F}&=\rho_{N,m}^{p,F}&&(m+1<N).
\end{align*}
It also leaves the free bed unchanged for $m<N$ and the entire
associated tower unchanged for $m+1<N$; its transport row is
unchanged already for $m<N$. These are identities of the maps and
coefficient data, not merely equalities of their evaluated actions.

For fixed $m$, all the identities hold once $N>m+1$. Deleting this
finite initial segment changes none of the limits above, none of
the eventual larger-width bounds and none of the retained free
horizon limits. Equality of the stationary maps gives equality of
every derivative at zero. Thus the final normalized construction
has exactly the same $R_m$, $\alpha_m$ and stationary jets. In
particular, replacing the terminal path has not introduced a second
background for the activity estimates.

\subsubsection{The auxiliary stationary service}

We supply the auxiliary input used in
Subsection~\ref{sec:selected-nonlinear-offering}. Fix one adjacent stage
pair, its paths and clips, its block side $L$, and a derivative order
$r$. Write $j=J_{m+1}$ for the incoming stage. All constants below may depend on these purchased data, but not
on the auxiliary period, forest or cutoff. The incoming loop allowance
is denoted by $\Lambda\ge1$. The auxiliary coarse side is at least two;
its fine side $Lc$ satisfies the larger-datum room bound
$Lc\ge100(j+7)L$. These are restrictions on $c$ after the stage
has been fixed. The selected kernel has the branch compatibility of
Subsubsection~\ref{subsec:path-stationary-chart}.

\begin{lemma}[Stationary matching on auxiliary periods]
\label{lem:selected-auxiliary-service}
The paths of the selected retained plan admit the uniform stationary
service on every sufficiently large auxiliary one-step torus. Its
critical graph agrees as a germ with the selected plan wherever that
plan is defined. On a coarse auxiliary side $c$, its classical matching
is exact and its thermal quadratic matching has bound
\begin{equation}\label{eq:selected-auxiliary-thermal}
 |D^2(\mathcal T_c-\mathcal D_c)(0)[h,h']|
 \le |\operatorname{Site}(c)|(A+\Gamma_j)\Lambda
              e^{-\min(\nu,c_j)c}
 \qquad(\|h\|_\infty,\|h'\|_\infty\le1),
\end{equation}
for the tested directions in the plan. Here $A<\infty$, $\nu>0$
are auxiliary service constants; $\Gamma_j,c_j>0$ are the retained
plan's thermal constants; and $\mathcal D_c$ is the draft formed from
its same outgoing path and scalar. The chart table is the plan's
local-ratio chart. These choices are uniform over the admissible
incoming paths with the purchased bounds.
\end{lemma}
\begin{proof}
\emph{Uniform local data.} Use the tree-gauge coordinates of
Subsubsection~\ref{subsec:path-stationary-chart}. A fixed block has a bounded
number $K$ of free coordinates per site, uniformly in the period and
allowed forest. Exponential path bounds give the local derivative
estimate \eqref{eq:stationary-scatter-input} at every fixed required
order. Here is the length payment in that assertion. Differentiating
a path word $r$ times selects at most $(1+\ell)^r$ positions. If the
selected coordinates are prescribed, the word visits all their anchors;
the sum of their distances from any one of them is at most $r\ell$.
From its weight $e^{\sigma\ell}$ reserve $e^{\sigma\ell/2}$ for
summation and use the rest to pay those distances and
$(1+\ell)^r$. Thus one can use a positive decay width smaller than
$\sigma/(4r)$ and a finite constant depending only on $r,\sigma$
and the path budget. Products of a bounded number of words obey the
same estimate after splitting the width between their factors.
The block coordinate maps, clip and local Haar logarithm have bounded
fixed-order derivatives on a fixed smaller coordinate ball; their
stencils have bounded diameter at this fixed $L$. Composition only
changes these constants and the decay width. Summing free derivative
indices costs the shell sum $K\sum_{n\ge0}(2n+1)^4e^{-an}$, independently
of the torus. This proves the local derivative input for the actual
path action. The classical action uses the fixed path allowance;
the loop derivatives have in addition their factor $\Lambda$.

The positive Hessian gap is the all-period tree-gauge gap of the same
selected horizon. Subsection~\ref{sec:uniform-stationary-chart} now
constructs a common coordinate ball and the unique critical graph
$\rho_c$, with positive smaller protection radii, smoothness through
the requested order, and a uniform inverse bound. Increase the local
order to at least $\max(r+2,4)$ when necessary. The argument there
uses supremum norms for the radii, so none shrink with the auxiliary
volume. Its estimates for loop and Haar derivatives give the remaining
chart bounds, with at most the stated loop factor.

We shall also need exponential entry bounds for the Hessian and its
first two derivatives along $\rho_c$. These follow from the same
local derivative estimate, not from a dimension-dependent norm
conversion. Indeed $D\rho=-H^{-1}D_xF$, and differentiating the
critical equation twice expresses $D^2\rho$ as $H^{-1}$ times a sum
of third derivatives of the action evaluated on first graph derivatives.
The inverse decays exponentially by Subsection~\ref{sec:local-covariance}.
Root each derivative at its unsummed matrix index and sum the remaining
indices with the preceding shell bound, reducing the decay width at
each of these finitely many convolutions. The chain rule for
$H(x,\rho_c(x))$ through order two then yields constants
$d,B,T,Q>0$, uniform in the period and forest, such that its inverse,
first jets and second jets have entry bounds respectively
$Be^{-d\operatorname{dist}}$, $Te^{-d\operatorname{dist}}$ and
$Qe^{-d\operatorname{dist}}$, for external directions of supremum norm
at most one. The gap fixes $B$; these constants involve classical
path derivatives only, so the loop factor is not squared.

\emph{Replication and uniqueness.} Choose a scheduled coarse side
$C=nc$, beyond the plan's volume threshold, using
Subsection~\ref{sec:common-scale-selection}. Lift the auxiliary forest
and its local coordinate conventions periodically. There are $n^4$
copies of each free coordinate. Replication copies the same path words,
block maps and local chart ratios. The stationary equation on a
replicated field is therefore the replication of the smaller stationary
equation. Both critical graphs tend to zero and belong to the common
uniqueness zone near zero. Their pullbacks under replication agree
there. The scheduled plan's graph has the same criticality and lies in
that zone after restricting its germ, so it too is the same graph.
This proves branch identification, rather than choosing a new branch
for the auxiliary estimates. Classical action and outgoing path sums
on replicated fields are exactly $n^4$ copies. The plan's classical
matching therefore descends after division by $n^4$.

\emph{The determinant error.} We give the finite-matrix calculation
because the thermal term does not replicate exactly. Let $\pi$ map
large free coordinates to their small-period addresses and let
$E_{ia}=\mathbf1_{\{\pi(i)=a\}}$. Write $S(x)$ for the large Hessian
on replicated coarse fields and $T(x)$ for the small Hessian. Equality
of the replicated stationary equations, differentiated in the free
variables, gives
\[
                    S(x)E=ET(x).
\]
Equivalently, the row sum of $S$ over a fibre equals the corresponding
entry of $T$. It is not an equality of individual inverse blocks.
Inverting and differentiating once or twice preserves this intertwining
relation, as does multiplication. If $PE=ER$, the equal fibre sizes
imply the exact trace identity
\begin{equation}\label{eq:selected-fibre-trace}
 n^4\operatorname{tr}R-\operatorname{tr}P
   =\sum_i\sum_{\substack{a\ne i\\\pi(a)=\pi(i)}}P_{ia}.
\end{equation}
Distinct coordinates in one fibre have anchor distance at least $c$.
With $S_K(t)=K\sum_{r\ge0}(2r+1)^4e^{-tr}$, an entry bound
$|P_{ia}|\le C_Pe^{-2t\operatorname{dist}(i,a)}$ consequently gives
\[
 |n^4\operatorname{tr}R-\operatorname{tr}P|
       \le |I_C|C_PS_K(t)e^{-tc}.
\]
This follows by extracting $e^{-tc}$ and summing the remaining
$e^{-t\operatorname{dist}}$ over all coordinates, which only enlarges
the sum over the fibre.

For clarity, exponential entry bounds are stable under the products
needed here: if two matrices have bounds $A e^{-a\operatorname{dist}}$
and $B'e^{-a\operatorname{dist}}$, then their product has bound
$AB'S_K(a/2)e^{-(a/2)\operatorname{dist}}$. Extract the last exponential
by the triangle inequality and sum a remaining exponential over the
intermediate index. Set $I=S(0)^{-1}$ and use
\[
 D^2\tfrac12\log\det S[h,h']
   =\tfrac12\operatorname{tr}(I S_{hh'}-I S_h I S_{h'}).
\]
The first product is bounded at width $d/4$ by
$BQ S_K(d/4)$, and each $I S_h$ is bounded at width $d/2$ by
$H_0=BT S_K(d/2)$. Thus the second product is bounded at width $d/4$
by $H_0^2 S_K(d/4)$. Apply \eqref{eq:selected-fibre-trace} with
$t=d/8$ to the two products. The absolute difference of the determinant
jets, with the small jet multiplied by $n^4$, is at most
\[
 |I_C|[BQ+H_0^2]S_K(d/4)S_K(d/8)e^{-dc/8}.
\]
The local Haar logarithm and the incoming loop sum replicate exactly
along the identified graphs, so their jets cancel in this difference.
There are at most $K L^4 C^4$ free coordinates. Dividing by $n^4$
therefore gives constants $A<\infty$, $\nu=d/8>0$ for which
\[
 |D^2\mathcal T_c-n^{-4}D^2(\mathcal T_C\circ E_c)|
        \le c^4 A e^{-\nu c},
\]
where $E_c$ is coarse-field replication and the directions have norm
at most one. This also explains the local-ratio requirement: the
Jacobian terms must be the same copied coordinate terms in both
expressions.

\emph{Transfer of the retained estimate.} Proposition
\ref{prop:path-retained-plan}, on the selected horizon and paths,
retains the scheduled estimate with amplitude $\Gamma_j\Lambda$ and
rate $c_j$. Its hypotheses are the same purchased one-step path
construction and quadratic horizon used at the beginning of this
section; no new stationary plan is selected for the auxiliary torus.
Apply that estimate to the replicated tested directions. They retain
supremum norm one, and the draft copies exactly. Division by $n^4$
turns its site factor into $c^4$. The triangle inequality yields
\[
 |D^2(\mathcal T_c-\mathcal D_c)|
 \le c^4\{A e^{-\nu c}+\Gamma_j\Lambda e^{-c_j C}\}.
\]
Since $C\ge c$ and $\Lambda\ge1$, this implies
\eqref{eq:selected-auxiliary-thermal}. All constants were fixed before
$c,C$ and the forest. The same reasoning applies to any allowed
auxiliary forest and its periodic lift; uniqueness identifies any
other admissible germ on it without changing the constants.
\end{proof}

\subsubsection{The activity and coupling limits on this background}
Fix a period $p$. The regular supports, forced inventories and norms
at scale $m$ depend on the fixed scale and period, not on $N$.
Embed their actual arrays into the one completion $\mathcal H_{m,p}$
of Subsection~\ref{sec:coupled-background-convergence}. The representative
of each finite history has exactly its prescribed activities; the
embedding is the same one used in the finite-test estimates. Its
norm is at most $1/8$. At every retained scale its path is
$\widetilde R_{N,m}$, by the construction's input/output identities.
The scalar in its coupling recursion is
$\widetilde\alpha_{N,m}$, not a newly chosen correction.

The finite real-test estimates and their ray consequence in
Subsection~\ref{sec:background-ray-estimates} give the two one-step
comparison bounds in
\eqref{eq:coupled-activity-step}--\eqref{eq:coupled-cost-step}.
Their constants are $8/128$ for the activity and $4s_m$ for the
scalar correction. The entry-threshold calculation at the end of
Subsection~\ref{sec:coupled-background-convergence} gives
$\sum_{j<m}s_j\le B_m/128$ for these same allowances. The
larger-width path limit above, the convergent stationary jets and
the retained free horizon supply the background data in the
common-variable comparison. The compensated parameter estimates
of Sections~\ref{sec:local-parameter-comparison} and
\ref{sec:finite-volume-parameter-continuity}, with the suppliers
of Subsection~\ref{sec:parameter-comparison-suppliers}, then give its
local continuity, including the projected scalar output. Each
application retains the funded transition and output identities
stated there. To apply this continuity to every auxiliary history
with a prescribed coupling script and zero terminal activity, one
must verify those transition hypotheses for its actual input at each
step. In particular, the protected chart, funding and renewed bounds,
downward-map bounds and output identification must refer to that input
and the same path and dictionaries. Availability of a one-step map on
the allowed ball alone does not establish these supplier hypotheses.
For the auxiliary zero-array inputs, Lemmas~\ref{lem:selected-auxiliary-service}
and~\ref{lem:selected-chartlet-matching} below verify this matching.
The ordered derivative bounds of the cited comparison argument are
not a further ingredient beyond this matching. Their supplier
construction in Subsection~\ref{sec:parameter-comparison-suppliers} is
stated at a fixed derivative order $j$ for the fine representation
satisfying the local hypotheses of Subsection~\ref{sec:activity-jet-assembly}.
It yields the jet bounds at every order through $j$ by the same local
estimates that give the value bounds: the elementary compensated
factors with their support-dependent derivative scale, the cluster jet
bounds, and the cluster criterion for $2^ju_P$ with its reserve. The
order needed in the continuity argument is at most the order $j$
supplied. Hence no value-order continuity theorem is required. What
remains is a matching statement, not a new estimate. The auxiliary
dictionary of Subsection~\ref{sec:selected-nonlinear-offering}, that is
the chartlet with zero regular activity at the stage coupling, must
satisfy the local hypotheses of
Sections~\ref{sec:activity-jet-assembly}
and~\ref{sec:parameter-comparison-suppliers}, on the same input, path
and dictionaries, with the healthy-mask, width and cutoff premises
retained there. Lemma~\ref{lem:selected-auxiliary-service} supplies its stationary service
on these same paths, and Lemma~\ref{lem:selected-chartlet-matching}
checks the remaining local hypotheses. Thus the auxiliary matching
required by the convergence argument is available on this construction.

\begin{lemma}[The auxiliary chartlet meets the comparison hypotheses]
\label{lem:selected-chartlet-matching}
Fix a stage and a point $x$ of an input interval. Consider the auxiliary
torus of Subsection~\ref{sec:selected-nonlinear-offering}, carrying the
same block bridge, the same two paths and clips, and the fine threshold,
radius and band at $\beta_F(x_0)$. Consider its reference fine
representation, the bare path action with the zero activity array.
Use the auxiliary stationary service of
Lemma~\ref{lem:selected-auxiliary-service}, on the same incoming path,
one-step datum and local-ratio chart. Then this
representation satisfies the incoming hypotheses of
Sections~\ref{sec:activity-jet-assembly}
and~\ref{sec:parameter-comparison-suppliers}, at every fixed derivative
order $j$, as follows:
\begin{enumerate}
\item \emph{Incoming activity budgets.} The regular and forced arrays
  are zero. Hence their regular and forced costs vanish, and every
  activity budget, packet hypothesis and cluster hypothesis imposed on
  them holds trivially.
\item \emph{Sector and shell access.} The block equation, the two
  adjacent kernel bounds and the all-period positive lower bounds are
  those of the selected horizon. They apply on every auxiliary torus, so
  the local expansion and its differentiated remainder bounds are
  available there. The shell-access constants are those of the same
  sector construction.
\item \emph{Warnings.} The deterministic safety and plateau inequalities
  of Subsection~\ref{sec:gaussian-warnings} hold, with their safety
  conditions produced in Subsection~\ref{sec:warning-safety}. That
  production concerns the fixed sector data and does not involve the
  activity array.
\item \emph{Protected chart.} The comparison chart is the stationary
  chart of Subsection~\ref{sec:uniform-stationary-chart}, restricted to the
  coarse horizon. Its identification with the graph used in the
  stationary estimates is the branch identification in Lemma~\ref{lem:selected-auxiliary-service}, on the same block, path and
  clipping data.
\item \emph{Running-curve data.} The renewed estimates are the zero
  costs of (1) together with the fine path bounds and the clock
  inequality for the loop allowance. The geometric bounds on the
  downward map are the kernel bounds of (2). The initial action is the
  admissible bare path action. The inequality
  $\operatorname{preExponent}/B\le\operatorname{width}/2$ is
  \eqref{eq:selected-offer-width}, which gives
  $A_c/B_c\le\sigma_c/4$. The clipping, the path and the regular and
  forced dictionaries are those of the physical comparison, with
  elementary labels identified by the identity. The displayed output is
  the constructed output of the local expansion on that torus.
\end{enumerate}
The healthy-support condition on a tested family of labels is not a
property of the representation. It is a premise of
Subsection~\ref{sec:activity-jet-assembly} for each family to which the
product estimate is applied, and it is retained there unchanged.
Consequently the supplier construction of
Subsection~\ref{sec:parameter-comparison-suppliers} applies to the
auxiliary chartlet at order $j$. It yields the ordered derivative
bounds at every order through $j$ by the same estimates as the value
bounds.
\end{lemma}

\begin{proof}
Each item is read off from the construction of the auxiliary data in
the subsection ``Keeping the physical and auxiliary data fixed'' of
Subsection~\ref{sec:selected-nonlinear-offering}, together with the
sections cited in the statement.

(1) The array is zero, so every coefficient vanishes. Regular and
forced costs are suprema of sums of nonnegative multiples of the
absolute values of these coefficients, hence zero. Every incoming
budget of the two cited sections is an upper bound by a nonnegative
allowance, so it holds.

(2) The cited subsection states that the auxiliary tori carry the same
block bridge and paths, and that the all-period kernel bounds apply on
every such torus, so that the local expansion and its differentiated
remainder bounds are available there. The shell-access constants are
constants of the sector construction on the selected horizon, which is
the same on these tori.

(3) The safety and plateau inequalities are deterministic statements
about the fixed sector, warning catalogue and thresholds, produced in
Subsection~\ref{sec:warning-safety}. They do not involve the activity
array.

(4) This is the protected-chart requirement of
Subsection~\ref{sec:parameter-comparison-suppliers}. That section
identifies a stationary chart with the graph of its estimates exactly
when its representative is a critical point of the same tested action
in the fine zone. On the auxiliary torus, this is the identification of
the auxiliary stationary branch with the plan's branch, proved in
Lemma~\ref{lem:selected-auxiliary-service}. Its eventual protected-field
condition holds because $\tau(x)\to0$.

(5) Each clause is quoted from the cited subsection and from
\eqref{eq:selected-offer-width}, since
$A_c/B_c\le\sigma_c/4\le\sigma_c/2$. The output identification holds
because the auxiliary output is, by definition, the output of the local
expansion applied to this representation.

Finally, Subsection~\ref{sec:parameter-comparison-suppliers} chooses its
thresholds before the incoming representation is supplied, and then
instantiates its constants at that representation. Its conclusions
therefore apply to every representation satisfying (1)--(5), and in
particular to the chartlet.
\end{proof}

With these same-input verifications, all hypotheses of
Subsection~\ref{sec:coupled-background-convergence} are on one and the
same family. For any terminal sequence
$N\to\infty$ and any corresponding finite trajectories whose bottom
couplings tend to $b$, there are, at every fixed $m,p$, limits
\[
 \beta_{N,m}\longrightarrow\beta_m,\qquad
 x_{N,m}\longrightarrow x_m\ \text{in }\mathcal H_{m,p},\qquad
 q_{N,m}\longrightarrow q_m.
\]
Here $q_{N,m}$ is the scalar activity correction, distinguished from
the path correction $\alpha_{N,m}$. They obey
\[
                \beta_m-\beta_{m+1}-\alpha_m=q_m.
\]
Two sequences with the same bottom limit $b$ give the same limits:
interleave them, whose terminal indices still tend to infinity,
and apply the same uniqueness proof. No compactness of a ball in
$\mathcal H_{m,p}$ is used. The norm convergence and the positive
partition lower bound in Subsection~\ref{sec:fixed-rung-convergence}
now apply to the actual normalized densities at that scale.

This application retains the strengthened stationary construction,
its convergent free horizon and the funded nonlinear transition.
The argument establishes compatibility of their consequences; the
analytic production of these inputs remains the responsibility of
the corresponding construction estimates.
